\documentclass[aspectratio=169]{beamer}

\usetheme{Madrid}
\usecolortheme{default}

\usepackage{amsmath}
\usepackage{graphicx}
\usepackage{hyperref}

\title{ECE 251: Computer Architecture}
\subtitle{Week 10: Multicycle Datapath \& Pipelining}
\author{Prof Rob Marano}
\date{Spring 2026}

% Cooper Union Brand Identity
\definecolor{CooperRed}{HTML}{8A1538}
\definecolor{CooperDark}{HTML}{5C0E25}
\setbeamercolor{structure}{fg=CooperRed}
\setbeamercolor{palette primary}{bg=CooperRed,fg=white}
\setbeamercolor{palette secondary}{bg=CooperRed,fg=white}
\setbeamercolor{palette tertiary}{bg=CooperDark,fg=white}
\setbeamercolor{title in head/foot}{bg=CooperDark,fg=white}
\logo{
\includegraphics[height=1cm]{../../images/cooper_union_logo.png}}
\titlegraphic{
\includegraphics[height=1.5cm,keepaspectratio]{../../images/cooper_union_logo.png}}

\begin{document}

\begin{frame}
\titlepage
\end{frame}

\begin{frame}{1. Single-Cycle Recap: The Bottleneck}
\begin{itemize}
    \item \textbf{Core Concept}: Every instruction completes entirely within 1 clock pulse. ($CPI = 1.0$).
    \item \textbf{The Reality}: Hardware propagation requires physical time (ps) to stabilize logic gates.
    \item \textbf{The Limitation}: The Clock Cycle Time ($T_c$) is determined by the absolute longest instruction format. For MIPS, the \texttt{lw} (Load Word) causes 800-1000ps delays.
    \item \textbf{Consequence}: The ALU \texttt{add} instruction inherently idles in compute time waiting for the global cycle to finish.
\end{itemize}
\end{frame}

\begin{frame}{2. The Multicycle Scheme}
\begin{itemize}
    \item \textbf{Solution}: Partition the Single-Cycle datapath into $3-5$ discrete executing phases.
    \item \textbf{Benefit 1}: Clock mapping is bounded by the slowest individual phase (e.g. Memory Access), permitting $200\text{ps}$ short clock periods.
    \item \textbf{Benefit 2}: Hardware units consolidate (e.g., sharing a single ALU instead of requiring dedicated PC Adders).
    \item \textbf{Cost}: Control structure transitions from simple logical assignment into a Finite State Machine (FSM).
\end{itemize}
\end{frame}

\begin{frame}{3. New Architectural Holding Registers}
Because data crosses clock boundaries in a Multicycle scheme, we implement holding registers between hardware components:
\vspace{0.2cm}
\begin{enumerate}
    \item \textbf{IR (Instruction Register)}: Latches the fetched instruction from Main Memory.
    \item \textbf{MDR (Memory Data Register)}: Holds data payload retrieved from RAM.
    \item \textbf{A \& B Registers}: Store operands extracted from the Register File.
    \item \textbf{ALUOut}: Captures the calculation output generated by the ALU.
\end{enumerate}
\end{frame}

\begin{frame}{4. Multicycle Universal States ($0-4$)}
Instructions are routed through the FSM:
\begin{itemize}
    \item \textbf{State 0 (Fetch)}: \texttt{Mem[PC] $\to$ IR}. Simultaneously, \texttt{PC + 4 $\to$ PC}.
    \item \textbf{State 1 (Decode)}: \texttt{Reg[rs] $\to$ A}, \texttt{Reg[rt] $\to$ B}. Calculate the Branch Target.
    \item \textbf{State 2 (Execute)}: ALU computes memory sum, logic functions, or branch conditions.
    \item \textbf{State 3 (Memory)}: Read or Write to the RAM array.
    \item \textbf{State 4 (Write Back)}: Register the finalized numerical calculation to the \texttt{rd} or \texttt{rt} destination.
\end{itemize}
\end{frame}

\begin{frame}[fragile]{5. SystemVerilog Emulation of the FSM}
\begin{verbatim}
module controller_fsm(input  logic       clk, reset,
                      input  logic [5:0] opcode,
                      output logic       memread, memwrite, regwrite);
                      
    typedef enum logic [3:0] {FETCH=0, DECODE=1, MEM_ADDR=2, MEM_READ=3, 
                              MEM_WRITE=4, R_EXEC=5, R_WRITEBACK=6} statetype;
    statetype state, next_state;
    
    // State Management
    always_ff @(posedge clk, posedge reset) 
        if(reset) state <= FETCH; else state <= next_state;
        
    // Combinational Next Logic based on Opcode
    always_comb begin
        case (state)
            FETCH:  next_state = DECODE;
            DECODE: case(opcode)
                       6'b100011: next_state = MEM_ADDR;
                       6'b000000: next_state = R_EXEC;
                       default:   next_state = FETCH;
                    endcase
            // ... additional branches
        endcase
    end
endmodule
\end{verbatim}
\end{frame}

\begin{frame}{6. The Multicycle Problem}
While the clock period latency is optimized, the hardware remains under-utilized.
\vspace{0.2cm}
\begin{itemize}
    \item During State 1 (Fetch), the mathematical ALU logic sleeps.
    \item During State 3 (Memory), the internal Register File is idle.
    \item The machine operates on one instruction traversing $3-5$ stages sequentially.
\end{itemize}
\end{frame}

\begin{frame}{7. The Concept of Pipelining (The Laundry Analogy)}
To extract system performance, we overlap the instruction execution phases.
\vspace{0.2cm}
\begin{center}
    \includegraphics[height=3.5cm]{"../../../images/textbook/ch004-9780128201091/jpg-9780128201091/004025.jpg"}
\end{center}
\vspace{0.2cm}
\begin{enumerate}
    \item If Wash takes 30 mins, Dry takes 30 mins, and Fold takes 30 mins, processing 4 laundry loads sequentially takes $4 \times 90 = 360$ mins.
    \item If we pipeline the sequence, the 4 loads complete in 180 mins.
\end{enumerate}
\end{frame}

\begin{frame}{8. The Reality of Pipelining: Hazards}
While pipelining overlaps execution to approach a $CPI = 1.0$, it physically forces instructions to share the datapath concurrently, creating timing conflicts and wait states called \textbf{Hazards}.
\vspace{0.2cm}
\begin{enumerate}
    \item \textbf{Structural Hazards}: The hardware lacks resources to support overlapping instructions (e.g., trying to read and write to a single memory array simultaneously).
    \item \textbf{Data Hazards}: An instruction depends on a data result from a preceding instruction that has not yet completed its Write Back phase.
    \item \textbf{Control Hazards}: The processor fetches sequential instructions, but a branch (\texttt{beq}) dynamically changes the Program Counter, rendering the fetched instructions invalid.
\end{enumerate}
\vspace{0.2cm}
When hazards occur, the pipeline must \textbf{stall} (wait), reducing performance until resolved via hardware bypassing.
\end{frame}

\begin{frame}{9. Week 10 Summary}
\begin{itemize}
    \item \textbf{Single-Cycle}: $CPI=1$, constrained by instruction length boundaries.
    \item \textbf{Multicycle}: FSM State execution, high clock speed, high $CPI$ mapping ($3-5$).
    \item \textbf{Pipeline}: Hardware overlap, aiming for $CPI = 1.0$ at high clock frequencies.
    \item \textbf{Next Session}: Analyzing and resolving Structural, Data, and Control Hazards.
\end{itemize}
\end{frame}

\end{document}
